\noindent\fbox{%
\begin{minipage}{\textwidth}
\textit{\textbf{KETENTUAN KHUSUS — OUTLINE JARINGAN KOMPUTER (Kode 137):}
\begin{enumerate}
    \item Tempat riset WAJIB sudah memiliki jaringan (walau sederhana).
    \item Surat Keterangan Riset: WAJIB ASLI — berkop surat, tertandatangani, berstempel basah. Harus ditunjukkan saat sidang.
    \item Skema Jaringan Berjalan (sub-bab 3.2.3) WAJIB di-capture dari software simulator: Packet Tracer, GNS3, Boson NetSim, atau sejenisnya — BUKAN digambar manual.
    \item Pengujian awal HARUS membuktikan permasalahan yang ada.
    \item Pengujian akhir HARUS membuktikan solusi usulan.
\end{enumerate}}
\end{minipage}%
}
\vspace{0.5cm}

\chapter{ANALISIS JARINGAN BERJALAN}
\section{Tinjauan Perusahaan}

\vspace{0.5cm}
\noindent\fbox{\begin{minipage}{\textwidth}
\textit{\textbf{Pedoman Penulisan:} Garis besar institusi/perusahaan tempat riset bergerak di bidang apa.}
\end{minipage}}
\vspace{0.5cm}

\subsection{Sejarah Perusahaan}
\subsection{Struktur Organisasi dan Fungsi}
\section{Skema Jaringan Berjalan}
\subsection{Topologi Jaringan}

\vspace{0.5cm}
\noindent\fbox{\begin{minipage}{\textwidth}
\textit{\textbf{Pedoman Penulisan:} Topologi yang digunakan di perusahaan tempat riset.}
\end{minipage}}
\vspace{0.5cm}

\subsection{Arsitektur Jaringan}

\vspace{0.5cm}
\noindent\fbox{\begin{minipage}{\textwidth}
\textit{\textbf{Pedoman Penulisan:} Arsitektur jaringan seperti IP Address, model OSI atau TCP/IP, QoS, dll.}
\end{minipage}}
\vspace{0.5cm}

\subsection{Skema Jaringan}

\vspace{0.5cm}
\noindent\fbox{\begin{minipage}{\textwidth}
\textit{\textbf{Pedoman Penulisan:} Informasi skema jaringan perusahaan (terdokumentasi maupun hasil penelitian di lapangan).}
\end{minipage}}
\vspace{0.5cm}

\subsection{Keamanan Jaringan}

\vspace{0.5cm}
\noindent\fbox{\begin{minipage}{\textwidth}
\textit{\textbf{Pedoman Penulisan:} Hal yang berhubungan dengan sistem keamanan seperti firewall, antivirus, proxy, manajemen bandwidth.}
\end{minipage}}
\vspace{0.5cm}

\subsection{Spesifikasi Hardware dan Software Jaringan}
\section{Permasalahan}

\vspace{0.5cm}
\noindent\fbox{\begin{minipage}{\textwidth}
\textit{\textbf{Pedoman Penulisan:} Menjabarkan masalah hardware maupun software yang ditemui.}
\end{minipage}}
\vspace{0.5cm}

\section{Alternatif Pemecahan Masalah}

\vspace{0.5cm}
\noindent\fbox{\begin{minipage}{\textwidth}
\textit{\textbf{Pedoman Penulisan:} Langkah-langkah untuk memecahkan masalah berdasarkan teori artikel ilmiah.}
\end{minipage}}
\vspace{0.5cm}

